% titel: Blatt 6, A6: Resonanzformeln von Stickstoffverbindungen

\subsection*{Lösung}
Für jede Verbindung die Resonanzstrukturen {Anzahl und Hauptmerkmale}:
\begin{itemize}
  \item \textbf{$\mathrm{N_2O}$} (lineares, 16 Valenzelektronen): zwei
        Resonanzformen        $\mathrm{N=N=O}$ und $\mathrm{N\equiv N-O}$; die
        reale Länge liegt zwischen $\mathrm{N-N}$- und $\mathrm{N=N}$-Bindung.
  \item \textbf{$\mathrm{NO}$} (Radikal, 11 $e^-$): zwei Resonanzstrukturen
        (eine mit $\mathrm{N=O}$, eine mit $\mathrm{N-O}$) plus ungepaartes
        Elektron.
  \item \textbf{$\mathrm{NO_2^+}$} (16 $e^-$): zwei äquivalente Strukturen mit
        je einer linearen $\mathrm{N=O}$-Doppelbindung und $\mathrm{N^+}$
        (linear).
  \item \textbf{$\mathrm{NO_2^-}$} (18 $e^-$): zwei Strukturen mit
        $\mathrm{N=O}$ / $\mathrm{N-O}$ und einer $-1$-Formalladung am
        freien O bzw.\ am N (gewinkelt, Resonanz stabilisiert beide).
  \item \textbf{$\mathrm{NO_3^-}$} (24 $e^-$): drei äquivalente
        Resonanzstrukturen, je eine $\mathrm{N=O}$-Doppelbindung und zwei
        $\mathrm{N-O}$-Einfachbindungen, Formalladung $-1$ auf dem einfach
        gebundenen O (trigonal-planar, Delokalisierung der Bindungslängen
        gleichwertig).
\end{itemize}

\begin{ergebnis}
$\mathrm{N_2O}$: 2, $\mathrm{NO}$: 2 (Radikal), $\mathrm{NO_2^+}$: 2,
$\mathrm{NO_2^-}$: 2, $\mathrm{NO_3^-}$: 3 äquivalente Resonanzstrukturen.
Alle werden zum realen Hybrid mittelwertet.
\end{ergebnis}

\begin{kontrolle}
Elektronenbilanz zur Sicherheit: $\mathrm{NO_3^-}$ hat $5+3\cdot6+1=24\,e^-$;
drei $\mathrm{N-O}$- (6) plus eine Doppelbindung (6) ergeben zusammen mit den
freien Paaren exakt 24. Die Zahl der Resonanzformen (3) stimmt mit der
Dreizähligkeit der Delokalisierung überein.
\end{kontrolle}

\begin{merksatz}
Resonanz = gleiche Atome, verschobene $\pi$-Bindungen und freie Paare;
Formalladungen ausgleichen, dann mittelt der Hybrid. Zähle die Elektronen, um
falsche Strukturen zu vermeiden.
\end{merksatz}
